\documentclass[11pt]{article}
\usepackage[T1]{fontenc}
\usepackage[utf8]{inputenc}
\usepackage{lmodern,microtype}
\usepackage[margin=1in]{geometry}
\usepackage{amsmath,amssymb,amsthm,mathtools}
\usepackage[hidelinks]{hyperref}
\usepackage{enumitem}
\newtheorem{theorem}{Theorem}[section]
\newtheorem{proposition}[theorem]{Proposition}
\newtheorem{lemma}[theorem]{Lemma}
\theoremstyle{remark}
\newtheorem{remark}[theorem]{Remark}
\newcommand{\R}{\mathbb R}
\newcommand{\loc}{\mathrm{loc}}
\newcommand{\dd}{\,dx}
\newcommand{\dA}{\,dA}
\newcommand{\ds}{\,ds}
\newcommand{\tr}{\operatorname{tr}}
\newcommand{\diver}{\operatorname{div}}
\newcommand{\supp}{\operatorname{supp}}
\newcommand{\eps}{\varepsilon}
\title{Second-order Sobolev regularity of three-dimensional
  \texorpdfstring{$p$}{p}-harmonic functions}
\author{Kaiwen Jin\\
  \small\href{mailto:yui.ui.siki@gmail.com}{yui.ui.siki@gmail.com}}
\date{October 13, 2026}
\begin{document}
\maketitle

\begin{abstract}
For every finite $p\geq5$, a weak $p$-harmonic function on an arbitrary
open subset of $\R^3$ belongs to $W^{2,2}_{\loc}$. The conclusion concerns
the full unweighted distributional Hessian and includes neighborhoods
of critical points without assumptions on the critical set.
We prove an estimate for smooth regularized solutions in which the
Hessian energy is divided by the actual gradient norm. Gaussian
curvature of the level surfaces supplies the component that is lost
in the usual pointwise coercivity estimate. The equation excludes
closed level components in a ball, so Gauss--Bonnet and the total
curvature of the boundary curves control the integrated correction.
A radial cutoff argument makes the estimate uniform in the
regularization parameter, and smooth approximation identifies every
second distributional derivative.
\end{abstract}

\section{Introduction}

The $p$-Laplace equation
\begin{equation}\label{eq:main}
 \diver\bigl(|Du|^{p-2}Du\bigr)=0
\end{equation}
is a basic model of degenerate ellipticity. Its weak solutions are
$p$-harmonic functions. Local H\"older continuity of the gradient is
available for every $1<p<\infty$, but for $p>2$ the coefficient in
\eqref{eq:main} degenerates where the gradient vanishes
\cite[Section 4]{Lindqvist}. Consequently, this first-order regularity
does not by itself identify a square-integrable Hessian across the
critical set.

Second-order estimates often control a nonlinear function of the
gradient. Sarsa proves local $W^{1,2}$ regularity of
$|Du|^{(p-2+\sigma)/2}Du$ when
$\sigma>-1-(p-1)/(n-1)$ \cite[Theorem 1.1]{Sarsa}.
Setting $\sigma=2-p$ removes the weight only in the range
$p<3+2/(n-2)$ for $n>2$. The same range appears in the classical
unweighted $W^{2,2}_{\loc}$ result discussed by Lindqvist
\cite[Section 11, p.~93, Second Derivatives]{Lindqvist} and in the global
Cordes-based theory of Haarala and Sarsa \cite{HaaralaSarsa}.
In three dimensions its upper endpoint is $5$. A weighted derivative
estimate does not settle what happens at a zero of $Du$ when the
weight vanishes, and the pointwise coercivity underlying that
argument does not remain positive beyond the endpoint.

We prove the unweighted conclusion in the remaining range. We use
the usual distributional definition: $u\in W^{1,p}_{\loc}(\Omega)$ is
a weak solution if
\begin{equation}\label{eq:weak}
 \int_\Omega |Du|^{p-2}Du\cdot D\varphi\dd=0
 \qquad\bigl(\varphi\in C_c^\infty(\Omega)\bigr).
\end{equation}

\begin{theorem}\label{thm:main}
Let $5\leq p<\infty$, let $\Omega\subset\R^3$ be open, and let
$u\in W^{1,p}_{\loc}(\Omega)$ be real valued and satisfy
\eqref{eq:weak}. Then $u\in W^{2,2}_{\loc}(\Omega)$.
In particular, each of the nine distributional derivatives
$D_jD_i u$, $1\leq i,j\leq3$, belongs to $L^2_{\loc}(\Omega)$.
\end{theorem}

Here $W^{k,r}_{\loc}$ means that the function and all its
distributional derivatives of order at most $k$ are locally in $L^r$.
Theorem~\ref{thm:main} imposes no sign or symmetry condition and no
restriction on the geometry or measure of the critical set.
It gives an affirmative answer to the three-dimensional
second-derivative question stated in Section~11 of Lindqvist's notes, including
$p=5$ and every larger finite exponent.

The missing normal component of the Hessian can be recovered by
adding signed Gaussian curvature of the level surfaces to a
divergence expression. This correction need not be positive at
each point. Its integral is nevertheless bounded by quantities on
the boundary of a ball: the equation rules out a closed regular
level component, and Gauss--Bonnet reduces its signed curvature
integral to the total curvature of the boundary curves. Fenchel's
inequality then controls the number of those curves.

The resulting estimate divides the squared Hessian by $|Dv|$ for
a smooth regularized solution $v$. At a critical point this
quotient requires care even though $v$ is smooth. A nonnegative
norm identity and Fatou's lemma establish its local integrability
and the favorable distributional inequality before curvature is
integrated. Radial cutoffs and an iteration remove the preliminary
dependence on the regularization parameter. The known uniform
gradient bound then converts this estimate into an unweighted
$L^2$ bound.

All constants below may depend on the fixed finite exponent $p$;
we make no assertion uniform as $p\to\infty$. The argument uses
two-dimensional level surfaces and does not assert the corresponding
result in higher spatial dimensions.
Section~\ref{sec:critical} treats critical points of smooth
regularizations. Section~\ref{sec:geometry} supplies the curvature
identity and its boundary estimate, and
Section~\ref{sec:energy} proves the uniform estimate and the
distributional conclusion. Appendix~\ref{sec:coarea} gives the
smooth critical-value and coarea facts used in the geometric step.
All vector norms are Euclidean and matrix norms are Hilbert--Schmidt;
balls in an argument have a common center.

\section{A norm identity through the critical set}\label{sec:critical}

Fix $5\leq p<\infty$, $\eps>0$, and a smooth real solution $v$ on an open set $U\subset\R^3$
of
\begin{equation}\label{eq:regularized}
 \diver\bigl((\eps+|Dv|^2)^{(p-2)/2}Dv\bigr)=0.
\end{equation}
The approximation to a general weak solution will be taken only
after the estimates for this fixed $v$ have been proved.
Write
\[
 q=|Dv|,\qquad H=D^2v,\qquad Z=\{q=0\}.
\]
On $U\setminus Z$, let $e=Dv/q$, and split the Hessian into its
normal and tangential blocks:
\begin{equation}\label{eq:blocks}
\begin{gathered}
 P=I-e\otimes e,\qquad a=e^THe,\qquad w=PHe,\\
 B=PHP\big|_{e^\perp},\qquad
 B_0=B-\tfrac12(\tr B)I_{e^\perp},\qquad
 s=1+(p-2)\frac{q^2}{q^2+\eps}.
\end{gathered}
\end{equation}
The block $B$ is a symmetric operator on the two-dimensional
tangent plane. Expansion of \eqref{eq:regularized} gives
\begin{equation}\label{eq:trace}
 \tr H+(s-1)a=0,\qquad \tr B=-sa,\qquad 1\leq s\leq p-1.
\end{equation}
In an orthonormal frame with first vector $e$,
\begin{equation}\label{eq:decomposition}
 |H|^2=(1+s^2/2)a^2+2|w|^2+|B_0|^2,\qquad
 T:=|H|^2-|He|^2=|w|^2+|B|^2.
\end{equation}
Since $s^2a^2\leq2|B|^2$ and $s\geq1$, these identities imply
\begin{equation}\label{eq:kato}
 |H|^2\leq3T.
\end{equation}

\begin{lemma}\label{lem:critical}
Define
\[
 E=\frac{|H|^2}{q}\quad\hbox{on }U\setminus Z,\qquad E=0\quad\hbox{on }Z.
\]
Then $E\in L^1_{\loc}(U)$. The weak gradient of $q$ is $He$,
with $e=0$ on $Z$. The vector field
\[
 G=(\Delta v)e,\qquad S=Dq-G
\]
has a $C^1$ extension for $G$ and a locally bounded realization for
$S$, both zero on $Z$. If
\begin{equation}\label{eq:f}
 f=\frac{|w|^2+|B_0|^2-\tfrac12s(s-2)a^2}{q}
 \quad\hbox{on }U\setminus Z,\qquad f=0\quad\hbox{on }Z,
\end{equation}
then $f\in L^1_{\loc}(U)$ and
\begin{equation}\label{eq:critical}
 \int_U \psi f\dd\leq-\int_U S\cdot D\psi\dd
 \qquad(0\leq\psi\in C_c^\infty(U)).
\end{equation}
Moreover, $|S|\leq(1+\sqrt3)|H|$ almost everywhere, and
\begin{equation}\label{eq:zeroidentity}
 |H|=\sqrt E\,\sqrt q\qquad\hbox{almost everywhere on }U.
\end{equation}
\end{lemma}

\begin{proof}
For $\delta>0$ set $q_\delta=(q^2+\delta^2)^{1/2}$. Differentiation gives
\[
 Dq_\delta=\frac{HDv}{q_\delta},\qquad
 \Delta q_\delta=
 \underbrace{\frac{|H|^2}{q_\delta}
   -\frac{|HDv|^2}{q_\delta^3}}_{T_\delta}
 +\frac{Dv}{q_\delta}\cdot D\Delta v.
\]
Here $T_\delta\geq0$ because $|HDv|\leq|H|q$.
On compact subsets, $q_\delta\to q$ uniformly and
$Dq_\delta\to He$ pointwise, including zero on $Z$;
also $|Dq_\delta|\leq|H|$. Dominated convergence in the test
identity for a first derivative identifies the weak gradient
$Dq=He$.
For a nonnegative compactly supported smooth $\psi$, integration by
parts yields
\[
 \int\psi T_\delta\dd
 =\int q_\delta\Delta\psi\dd
   -\int\psi\frac{Dv}{q_\delta}\cdot D\Delta v\dd.
\]
The right side has a finite limit: $D\Delta v$ is bounded on the
support and $|Dv/q_\delta|\leq1$. Fatou's lemma, applied on $U\setminus Z$,
therefore gives
\begin{equation}\label{eq:fatou}
 \int_{U\setminus Z}\psi\frac Tq\dd
 \leq-\int Dq\cdot D\psi\dd
       -\int\psi e\cdot D\Delta v\dd.
\end{equation}
Choose $\psi$ equal to one on a prescribed compact set.
Equations~\eqref{eq:kato} and \eqref{eq:fatou} imply
$E\in L^1_{\loc}$. This is a finiteness statement for the
fixed smooth solution; its present bound may depend on $\eps$
and on third derivatives.

We next justify integration by parts for $G$ at $Z$.
Equation~\eqref{eq:trace} implies
\begin{equation}\label{eq:G}
 G=-(p-2)\frac{(Dv^THDv)Dv}{q(q^2+\eps)}\qquad(q>0).
\end{equation}
For a symmetric matrix $A$, the map
$F(z,A)=(z^TAz)z/|z|$, assigned zero at $z=0$, is $C^1$.
Indeed, off zero
\[
 |F|\leq |A||z|^2,\qquad
 |D_zF|\leq C|A||z|,\qquad |D_AF|\leq C|z|^2.
\]
The first bound shows that the derivative at zero is zero,
and the other bounds give its continuity there. Composing with
$Dv,H$ and dividing by the positive smooth function $q^2+\eps$
in \eqref{eq:G} proves that $G$ is $C^1$ and $G=DG=0$ on $Z$.

Off $Z$, differentiation of $e$ gives
$\diver e=(\Delta v-a)/q$, so the classical divergence of $S$ is
\[
 \frac{T-(\Delta v)^2+(\Delta v)a}{q}.
\]
Using $\Delta v=(1-s)a$ and \eqref{eq:decomposition} gives exactly
\eqref{eq:f}. Equivalently, off $Z$,
\[
 f=T/q+e\cdot D\Delta v-\diver G.
\]
All three terms are locally integrable. The latter two vanish on
$Z$, and $T/q$ is extended by zero there. Adding those two terms
to \eqref{eq:fatou} and integrating $G$ by parts proves
\eqref{eq:critical}. In particular, the passage through the
critical set preserves the required direction of the inequality.
Since $|Dq|\leq|H|$ and $|\Delta v|\leq\sqrt3|H|$, the stated bound
for $S$ follows.

Finally $H=0$ almost everywhere on $Z$. For each smooth component
$v_i$ of $Dv$, the part of $\{v_i=0\}$ where $Dv_i\ne0$ is locally
a smooth hypersurface by the implicit function theorem. A graph
has three-dimensional measure zero by Fubini, and countably many
such graphs cover that part of the zero set. Apply this to all
three components $v_i$. At almost every point of $Z$ all the
rows of $H$ vanish, proving \eqref{eq:zeroidentity}.
The convergence and integration facts used here are the usual
Fatou, dominated convergence and Fubini theorems
\cite[Sections 1.4 and 1.7]{TaoMeasure}; the smooth calculus and
implicit graph facts can be found in
\cite[Sections 4.3 and 5.3]{LeblI} and
\cite[Sections 8.3, 8.5 and 8.6]{LeblII}.
\end{proof}

\section{Curvature of level surfaces}\label{sec:geometry}

\subsection{The coercive identity}

At $q>0$, the second fundamental form of the level surface of $v$
is $B/q$, up to the sign convention for the normal. Its Gaussian
curvature is therefore $K=\det B/q^2$. Introduce the signed density
\begin{equation}\label{eq:curvature}
 \mathcal K=Kq=\frac{\det B}{q}
   =\frac{s^2a^2/4-|B_0|^2/2}{q}
 \quad(q>0),\qquad \mathcal K=0\quad(q=0).
\end{equation}
The determinant of a symmetric two-dimensional matrix satisfies
$|\det B|\leq|B|^2/2$, whence
$|\mathcal K|\leq E/2$.
Lemma~\ref{lem:critical} has thus established absolute integrability
before we integrate signed curvature.

Let
\begin{equation}\label{eq:constants}
 \ell=p-1,\qquad c_p=\ell+\ell^{-1},\qquad d_p=2c_p-2.
\end{equation}
Substitution of \eqref{eq:decomposition}, \eqref{eq:f}, and
\eqref{eq:curvature} gives the identity
\begin{equation}\label{eq:coerciveidentity}
 c_pf+d_p\mathcal K-E
 =\frac{(c_p-2)|w|^2+(sc_p-s^2-1)a^2}{q}\qquad(q>0).
\end{equation}
The coefficient of $|B_0|^2$ cancels. The remaining normal
coefficient factors as
\[
 sc_p-s^2-1=\frac{(\ell-s)(s\ell-1)}{\ell}\geq0
 \qquad(1\leq s\leq\ell),
\]
and $c_p-2=(\ell-1)^2/\ell\geq0$. Consequently
\begin{equation}\label{eq:coercive}
 E\leq c_pf+d_p\mathcal K
 \qquad\hbox{almost everywhere on }U.
\end{equation}
There is no separate positivity assertion for $f$ or for $K$.

\subsection{An integrated boundary estimate}

\begin{lemma}\label{lem:boundary}
For every closed ball $\overline{B_r}\subset U$,
\begin{equation}\label{eq:boundary}
 \int_{B_r}\mathcal K\dd
 \leq2\int_{\partial B_r}\left(|H|+\frac{2q}{r}\right)\dA.
\end{equation}
\end{lemma}

\begin{proof}
Translate the center to zero. By Lemma~\ref{lem:values}, almost
every $t$ is a regular value both of $v$ near the ball and of
$v|_{\partial B_r}$. For such $t$,
$\Sigma_t=\{v=t\}\cap\overline{B_r}$ is a compact smooth
orientable surface with smooth boundary, possibly empty.
Its orientation comes from $Dv/q$.
It has finitely many components, as does its boundary: connected
coordinate neighborhoods have a finite subcover of a compact
manifold and each belongs to one component.

Every nonempty component has boundary. Otherwise a connected
closed embedded component $M$ lies in $B_r$. It separates a bounded
region $D$ from the unbounded region
\cite[B19]{SullivanSurfaces}. Every point outside the ball can
be joined to infinity without meeting $M$, so $D\subset B_r$,
and $v=t$ on $\partial D=M$.
To use the equation without importing a boundary regularity
argument, freeze the coefficient field
\[
 A(x)=I+(p-2)\frac{Dv\otimes Dv}{q^2+\eps},\qquad
 Lz=\tr(A(x)D^2z).
\]
Then $Lv=0$ and $I\leq A\leq(p-1)I$. For $\rho>0$,
$L(v+\rho e^{x_1})=\rho A_{11}e^{x_1}>0$.
At an interior maximum the Hessian is negative semidefinite,
which would give $L(v+\rho e^{x_1})\leq0$.
Thus its maximum on $\overline D$ is attained on $\partial D$.
Letting $\rho\downarrow0$ gives $v\leq t$ in $D$.
Applying the minimum argument to $v-\rho e^{x_1}$ gives
$v\geq t$. Hence $v=t$ in $D$, and continuity of $Dv$ gives
$Dv=0$ on $M$, contradicting regularity of the level.

Let $b$ be the total number of boundary curves. A connected
orientable compact surface of genus $\gamma_i$ with $b_i$
boundary circles has Euler characteristic $2-2\gamma_i-b_i$
\cite[B17]{SullivanSurfaces}. Here $b_i\geq1$, so
$\chi(\Sigma_t)\leq\sum_i b_i=b$.
Gauss--Bonnet for a smooth boundary
\cite[B16--B18]{SullivanSurfaces} gives
\begin{align}
 \int_{\Sigma_t}K\dA
 &=2\pi\chi(\Sigma_t)-\sum_\Gamma\int_\Gamma k_g^{\Sigma_t}\ds
 \notag\\
 &\leq2\pi b+\sum_\Gamma\int_\Gamma|k_g^{\Sigma_t}|\ds
 \leq2\sum_\Gamma\int_\Gamma\kappa\ds. \label{eq:GB}
\end{align}
Here $k_g^{\Sigma_t}$ is signed geodesic curvature with the
Gauss--Bonnet orientation, and $\kappa$ is nonnegative Euclidean
curvature of a boundary curve. The last inequality uses
$|k_g^{\Sigma_t}|\leq\kappa$ and Fenchel's inequality
$2\pi\leq\int_\Gamma\kappa\ds$ separately for each closed
regular curve \cite[Theorem 2.4 and Proposition 3.1]{SullivanCurves}.

We bound this last curvature on the sphere. Put $n=x/r$ and
$g=|D_{\partial B_r}v|$. On $\{g>0\}$ an arclength parametrization
$\Gamma$ of a spherical level curve has unit tangent $T$ and
\[
 \Gamma''=\eta\nu-\frac nr,\qquad
 \nu=\frac{D_{\partial B_r}v}{g},
\]
for a signed scalar $\eta$. Indeed, differentiation of $n\cdot T=0$
gives the normal component $-n/r$, while the tangential component
is perpendicular to $T$.
Differentiating $v(\Gamma)=t$ twice gives
$0=H(T,T)+\eta g-D_nv/r$. It follows that
\begin{equation}\label{eq:spherecurvature}
 \kappa g\leq|\eta|g+\frac gr
 \leq |H|+\frac{|D_nv|+g}{r}
 \leq |H|+\frac{2q}{r}.
\end{equation}
The volume coarea identity in Lemma~\ref{lem:coarea} gives
\[
 \int_{B_r}\mathcal K\dd
   =\int_\R\int_{\Sigma_t}K\dA\,dt.
\]
It applies to this signed integrand because $|\mathcal K|\leq E/2$
is integrable. Surface coarea on $\partial B_r$ gives
\[
 \int_\R\sum_\Gamma\int_\Gamma\kappa\ds\,dt
   =\int_{\partial B_r}\kappa g\dA,
\]
with $\kappa g=0$ on $\{g=0\}$. This nonnegative integral is finite
by \eqref{eq:spherecurvature}. Integrate \eqref{eq:GB} in $t$
and use \eqref{eq:spherecurvature} to obtain \eqref{eq:boundary}.
The critical-value assertion concerns the fixed smooth $v$,
and does not make any assertion about critical values of the
limiting weak solution.
\end{proof}

\section{Uniform energy and distributional second derivatives}\label{sec:energy}

\begin{proposition}\label{prop:energy}
For each finite $p\geq5$ there is $L_p<\infty$, depending only
on $p$, such that every smooth solution of
\eqref{eq:regularized} on a neighborhood of $\overline{B_{2R}}$
satisfies
\begin{equation}\label{eq:energy}
 \int_{B_R}\frac{|D^2v|^2}{|Dv|}
       \mathbf1_{\{|Dv|>0\}}\dd
 \leq\frac{L_p}{R^2}\int_{B_{2R}}|Dv|\dd.
\end{equation}
The constant is independent of $\eps>0$ and of $v$.
\end{proposition}

\begin{proof}
Retain $q=|Dv|$, $H=D^2v$, and the zero-extended densities
$E,\mathcal K$ from Sections~\ref{sec:critical} and
\ref{sec:geometry}. Set
$F(r)=\int_{B_r}E\dd$ and $Q=\int_{B_{2R}}q\dd$.
These are finite for this fixed smooth solution.
For $R\leq a<b\leq2R$ put $h=b-a$, and choose a smooth
decreasing function $\eta$, equal to one for $r\leq a$ and
zero for $r\geq b$, with $0\leq\eta\leq1$ and
$|\eta'|\leq2/h$. Such a function is obtained by integrating
a smooth nonnegative unit-mass function supported in $(0,1)$
and bounded by $2$, then scaling the interval. The radial
function $\psi(x)=\eta(|x|)$ is a smooth compactly supported
test, constant near the center.
Equations~\eqref{eq:critical} and \eqref{eq:coercive} yield
\begin{equation}\label{eq:testenergy}
 F(a)\leq\int\psi E\dd
 \leq-c_p\int S\cdot D\psi\dd+d_p\int\psi\mathcal K\dd.
\end{equation}

The identity $\psi(x)=\int_{|x|}^b-\eta'(r)\,dr$ and Fubini give
\[
 \int\psi\mathcal K\dd
 =\int_a^b-\eta'(r)\int_{B_r}\mathcal K\dd\,dr.
\]
The interchange is justified by absolute integrability of
$\mathcal K$, despite its sign. Apply Lemma~\ref{lem:boundary}
under this nonnegative radial weight, and use polar integration
on the shell $B_b\setminus B_a$. The result is
\begin{equation}\label{eq:radialcurvature}
 \int\psi\mathcal K\dd
 \leq\frac4h\int_{B_b\setminus B_a}|H|\dd
       +\frac8{ah}\int_{B_b\setminus B_a}q\dd.
\end{equation}
Polar integration here is the smooth coarea identity for $|x|$
on the shell.
Also, Lemma~\ref{lem:critical} gives
\[
 -\int S\cdot D\psi\dd
 \leq\frac{2(1+\sqrt3)}h\int_{B_b\setminus B_a}|H|\dd.
\]
Put
\[
 A_p=2c_p(1+\sqrt3)+4d_p,\qquad B_p=8d_p.
\]
By \eqref{eq:zeroidentity}, Cauchy--Schwarz, and
\eqref{eq:testenergy}--\eqref{eq:radialcurvature},
\begin{align}
 F(a)&\leq\frac{A_p}{h}\sqrt{F(b)Q}+\frac{B_p}{ah}Q\notag\\
 &\leq\vartheta F(b)
   +\frac{A_p^2/(4\vartheta)+B_p}{h^2}Q
   \qquad(0<\vartheta<1). \label{eq:iteration}
\end{align}
The second line uses $xy\leq\vartheta x^2+y^2/(4\vartheta)$
and $h\leq R\leq a$.

Choose $\vartheta=1/8$ and write
$L=A_p^2/(4\vartheta)+B_p$.
With $r_j=2R-R2^{-j}$, iteration of \eqref{eq:iteration} gives
\[
 F(R)\leq\vartheta^mF(r_m)
    +\frac{4LQ}{R^2}\sum_{j=0}^{m-1}(4\vartheta)^j.
\]
For the fixed $\eps$ and $v$, $F(r_m)\leq F(2R)<\infty$.
Thus the terminal term tends to zero as $m\to\infty$.
Since $4\vartheta=1/2$, \eqref{eq:energy} follows with $L_p=8L$.
This iteration is completed before taking any limit in $\eps$;
no uniform bound on the preliminary quantity $F(2R)$ was used.
\end{proof}

\begin{proof}[Proof of Theorem~\ref{thm:main}]
Fix an arbitrary ball $\overline{B_{4R}}\subset\Omega$.
The local $C^{1,\alpha}$ representative of $u$ is available
by \cite[Section 4]{Lindqvist}. Use the smooth Dirichlet
approximation on $U=B_{4R}$ described in
\cite[Section 3, (3.2)--(3.4)]{Sarsa}: for a sequence
$\eps_k\downarrow0$, there are $v_k\in C^\infty(U)$ solving
\eqref{eq:regularized} with $\eps=\eps_k$ and boundary data $u$,
such that the interior $C^{1,\alpha}$ norms are uniformly
bounded and $Dv_k\to Du$ locally uniformly.
These approximation statements hold for every finite $p>1$;
they do not require a second-order bound on $u$.
In the present application, $u$ is bounded on $\overline U$
by its local $C^1$ regularity, the ball is smooth, and
$\overline{B_{2R}}\subset U$ has positive distance from its boundary.
Consequently
\[
 M:=\sup_k\sup_{B_{2R}}|Dv_k|<\infty.
\]

Lemma~\ref{lem:critical} says that $D^2v_k=0$ almost everywhere
where $Dv_k=0$. Proposition~\ref{prop:energy} therefore gives
the full unweighted bound
\begin{equation}\label{eq:unweighted}
 \int_{B_R}|D^2v_k|^2\dd
 \leq\frac{L_pM}{R^2}\int_{B_{2R}}|Dv_k|\dd
 \leq\frac{L_pM^2|B_{2R}|}{R^2}=:C_0^2.
\end{equation}
This includes the case $M=0$.
For every $i,j\in\{1,2,3\}$ and
$\varphi\in C_c^\infty(B_R)$, local uniform gradient convergence
and integration by parts imply
\[
 -\int_{B_R}D_i u\,D_j\varphi\dd
   =\lim_{k\to\infty}\int_{B_R}D_jD_i v_k\,\varphi\dd.
\]
By \eqref{eq:unweighted}, its absolute value is at most
$C_0\|\varphi\|_{L^2(B_R)}$. Thus $D_jD_i u$, initially a
distribution, defines a bounded linear functional on
$C_c^\infty(B_R)$ with the $L^2$ norm.

For completeness, density on this ball follows from the usual
$C_c^\infty(\R^3)$ density in $L^2(\R^3)$ and smooth cutoffs
\cite[Exercise 1.13.1(iii)--(iv)]{TaoEpsilon}.
Extend $h\in L^2(B_R)$ by zero and take cutoffs
$0\leq\eta_m\leq1$ supported in $B_R$ with
$\eta_m\to1$ pointwise there. Dominated convergence gives
$\eta_mh\to h$ in $L^2$. For fixed $m$, approximate the zero
extension of $h$ by smooth compactly supported functions on
$\R^3$ and multiply by $\eta_m$. A diagonal choice proves
density of $C_c^\infty(B_R)$ in $L^2(B_R)$.
The functional extends uniquely to that real Hilbert space.
Riesz representation \cite[Theorem 1.4.13]{TaoEpsilon} supplies
$h_{ji}\in L^2(B_R)$ such that
\[
 -\int_{B_R}D_i u\,D_j\varphi\dd
   =\int_{B_R}h_{ji}\varphi\dd
   \qquad(\varphi\in C_c^\infty(B_R)).
\]
This is precisely the distributional identity $D_jD_i u=h_{ji}$.
It treats all nine entries without any condition on the
limiting critical set. Because $p\geq5>2$ and $B_R$ has finite
measure, $u,Du\in L^2(B_R)$ by H\"older's inequality.
Hence $u\in W^{2,2}(B_R)$.
Every point of $\Omega$ admits such balls. Distributional
uniqueness identifies the derivatives on overlaps, and finite
covers of compact subsets give their local square integrability.
\end{proof}

\appendix
\section{Smooth critical values and coarea}\label{sec:coarea}

The geometric argument needs these facts only for a fixed smooth
regularization and its restriction to a sphere. The following
proofs also specify the measures in the two coarea identities.

\begin{lemma}\label{lem:values}
If $f$ is smooth on an open subset of $\R^m$, $1\leq m\leq3$,
then the set of values at points where $Df=0$ has
one-dimensional Lebesgue measure zero. The same holds on a smooth
$m$-dimensional manifold with a countable atlas.
\end{lemma}

\begin{proof}
Induct on $m$. Let $A_j$ be the set where all derivatives of
orders $1,\ldots,j$ vanish. On a compact subset of the domain,
divide a bounding rectangle into cubes of side $h$.
For a cube meeting $A_m$, choose $x$ in that intersection.
Taylor's theorem applied to $\tau\mapsto f(x+\tau(y-x))$
gives $|f(y)-f(x)|\leq Ch^{m+1}$ for $y$ in that cube:
the derivatives through order $m$ vanish at $x$, and the
$(m+1)$st derivatives are bounded on a fixed neighborhood.
There are at most $C'h^{-m}$ cubes. The outer measure of the
image of $A_m$ on this compact subset is therefore at most
$C''h$, and hence zero. A compact exhaustion removes the
compactness restriction.

For $1\leq j<m$ and $x\in A_j\setminus A_{j+1}$,
some derivative $D^\beta f$ of order $j$ has a nonzero
first derivative at $x$. Its zero set is a smooth
$(m-1)$-dimensional hypersurface near $x$ and contains $A_j$
there. Since $Df=0$ on $A_j$, these points are critical
for the restriction of $f$ to that hypersurface. The induction
hypothesis, in its charts, makes their image null.
Countably many such neighborhoods cover each stratum.
Their union with $A_m$ is $A_1$. For $m=1$ only the Taylor
argument is needed. Countably many charts also establish the
manifold statement.
Taylor's theorem and the implicit function theorem used here
are \cite[Theorem 4.3.2]{LeblI} and
\cite[Theorem 8.5.6]{LeblII}.
\end{proof}

\begin{lemma}\label{lem:coarea}
Let $f$ be smooth on an open subset $V$ of $\R^3$. For a Borel
function $h$ on $\{|Df|>0\}$,
\begin{equation}\label{eq:volumecoarea}
 \int_{\{|Df|>0\}}h|Df|\dd
   =\int_\R\int_{\{f=t\}\cap\{|Df|>0\}}h\dA\,dt
\end{equation}
if $h\geq0$, or if the integral of $|h||Df|$ is finite.
For a smooth surface $N$ and $g=|D_Nf|$, the corresponding identity is
\begin{equation}\label{eq:surfacecoarea}
 \int_{\{g>0\}\subset N}hg\dA
   =\int_\R\int_{\{f=t\}\cap\{g>0\}}h\ds\,dt
\end{equation}
under the analogous condition.
\end{lemma}

\begin{proof}
In a coordinate rectangle with $f_1\ne0$, solve $f(x_1,y)=t$
for $x_1$. One-variable substitution gives
$dx_1=dt/|f_1|$, and the level graph has area element
$|Df|\,dy/|f_1|$. These equalities and Fubini give
\eqref{eq:volumecoarea} on the rectangle.
To justify this substitution for Borel integrands, first
apply the elementary substitution theorem to intervals
\cite[Theorem 5.3.5]{LeblI}. The two finite measures agree
on their Boolean algebra of finite interval unions and
therefore on Borel sets by the monotone-class lemma
\cite[Lemma 1.7.14]{TaoMeasure}. Exhaustion, simple functions
and monotone convergence give the nonnegative formula.
Positive and negative parts give the absolutely integrable
formula.

For the surface identity, write a chart $X(y_1,y_2)$ with
metric $G_{ij}=X_i\cdot X_j$. On a chart where $f_1\ne0$,
the speed of the level curve parametrized by $y_2$ is
\[
 L=\left|X_2-\frac{f_2}{f_1}X_1\right|,\qquad
 g^2=\frac{G_{22}f_1^2-2G_{12}f_1f_2+G_{11}f_2^2}{\det G}.
\]
Thus $L=g\sqrt{\det G}/|f_1|$, while substitution of $t$
in the surface area element gives
$\sqrt{\det G}\,dt\,dy_2/|f_1|$. This proves
\eqref{eq:surfacecoarea} in the chart by the same substitution
argument. Cover each noncritical set with countably many
such charts, partition it into disjoint Borel subsets of
the charts and sum using Tonelli. Absolute integrability
permits the signed sum. Lemma~\ref{lem:values} allows the
$t$ integrals to be restricted to simultaneous regular values.
\end{proof}

\section*{AI-use statement}
AI assistance was used in developing and checking the argument,
locating and checking sources, and preparing the manuscript.
The submission is offered for independent mathematical review.

\begin{thebibliography}{9}
\bibitem{Lindqvist}
P. Lindqvist,
\emph{Notes on the $p$-Laplace equation}, second edition,
Report 161, University of Jyv\"askyl\"a, 2017.
\href{https://jyx.jyu.fi/bitstreams/2a8cbc94-58e9-4e75-94d6-0f1f5bc22ac1/download}
{Author's monograph}.

\bibitem{Sarsa}
S. Sarsa,
\emph{Note on an elementary inequality and its application to the
regularity of $p$-harmonic functions},
arXiv:2009.10102v1, 2020.
\url{https://arxiv.org/abs/2009.10102v1}.

\bibitem{HaaralaSarsa}
A. Haarala and S. Sarsa,
\emph{Global second order Sobolev-regularity of $p$-harmonic functions},
Journal of Differential Equations \textbf{339} (2022), 232--260.
\href{https://doi.org/10.1016/j.jde.2022.08.022}
{doi:10.1016/j.jde.2022.08.022}.
Preprint: \url{https://arxiv.org/abs/2204.13550v2}.

\bibitem{SullivanSurfaces}
J. M. Sullivan,
\emph{Diff Geom I, SS 2013}, section B: Surfaces,
TU Berlin course notes, summer semester 2013, B16--B19.
\url{https://www3.math.tu-berlin.de/geometrie/Lehre/SS13/DGI/dg1.pdf}.

\bibitem{SullivanCurves}
J. M. Sullivan,
\emph{Curves of Finite Total Curvature},
arXiv:math/0606007v2, 24 October 2007.
\url{https://arxiv.org/abs/math/0606007v2}.

\bibitem{TaoMeasure}
T. Tao,
\emph{An Introduction to Measure Theory},
American Mathematical Society, 2011.
Author's preliminary version:
\url{https://terrytao.wordpress.com/wp-content/uploads/2012/12/gsm-126-tao5-measure-book.pdf}.

\bibitem{TaoEpsilon}
T. Tao,
\emph{An epsilon of room: pages from year three of a mathematical blog},
author's preliminary version, 2010, Chapter 1.
\url{https://terrytao.wordpress.com/wp-content/uploads/2010/02/epsilon.pdf}.

\bibitem{LeblI}
J. Lebl,
\emph{Basic Analysis I: Introduction to Real Analysis, Volume I},
version 6.3, 15 May 2026.
\url{https://www.jirka.org/ra/realanal.pdf}.

\bibitem{LeblII}
J. Lebl,
\emph{Basic Analysis II: Introduction to Real Analysis, Volume II},
version 6.3, 15 May 2026.
\url{https://www.jirka.org/ra/realanal2.pdf}.
\end{thebibliography}
\end{document}
